\documentclass[12pt]{article}

\usepackage{graphicx}
\usepackage{booktabs}

\title{Replication of eBay Paid Search Experiment}
\author{Patrick Michael}
\date{\today}

\begin{document}

\maketitle

\section{Introduction}

This paper replicates key findings from the eBay paid search experiment. 
The goal is to examine the causal impact of paid search advertising on firm revenue. 
Using experimental data, we compare treated and untreated groups over time and estimate the difference-in-differences effect of advertising.

\section{Experimental Design}

The experiment randomly assigned users to treatment and control groups. 
The treatment group was exposed to paid search advertising while the control group was not. 
Revenue outcomes were then observed over multiple time periods.

Because assignment was randomized, the treatment effect can be estimated using a difference-in-differences framework comparing the average revenue across groups and time.

\section{Revenue Trends}

Figure \ref{fig:revtime} shows the average revenue for treated and untreated groups over time.

\begin{figure}[h]
\centering
\includegraphics[width=0.8\textwidth]{../output/figures/figure_5_2.png}
\caption{Average Revenue by Group Over Time}
\label{fig:revtime}
\end{figure}

Figure \ref{fig:loggap} shows the difference in log revenue between the treated and untreated groups.

\begin{figure}[h]
\centering
\includegraphics[width=0.8\textwidth]{../output/figures/figure_5_3.png}
\caption{Difference in Log Revenue Between Treated and Control}
\label{fig:loggap}
\end{figure}

The figures suggest that the treated group experiences higher revenue during the experimental period.

\section{Difference-in-Differences Results}

Table \ref{tab:did} presents the estimated difference-in-differences effect of paid search advertising.

\input{../output/tables/did_table.tex}

The estimate represents the causal impact of advertising on revenue. 
A positive estimate would suggest that paid search advertising increases firm revenue.

\section{Conclusion}

This replication exercise demonstrates how experimental data can be used to estimate causal effects. 
The results illustrate the value of randomized experiments in measuring the effectiveness of advertising campaigns.

\end{document}
